\section{Conclusion \& Future Research Directions}
In this paper, we presented \textbf{C-STGB} (\underline{C}onformal \underline{S}patio-\underline{T}emporal \underline{G}raph\underline{B}oost), an autonomous, geometric, and physics-informed framework designed for robust anti-money laundering surveillance under extreme class skew ($\pi < 0.05\%$), continuous-time velocity evasion, and adversarial merchant camouflage. 

By formulating a continuous-time multi-scale temporal attention kernel with dynamic burst gating, C-STGB resolves both microsecond smurfing bursts and 90-day hibernation holding strategies. The learnable edge-gated topological filter suppresses over 65.9\% of adversarial camouflage links, eliminating the over-smoothing failure mode characteristic of spatial GNNs. Within the latent manifold, Latent-Space GraphSMOTE coupled with a parametric bilinear link generator synthesizes realistic illicit topology, effectively recovering standalone GNN recall from 10.33\% to 100.00\%. Furthermore, integrating Class-Conditional Conformal Risk Control (CRC) guarantees finite-sample error bounds ($\mathbb{P}(Y \in \Gamma(X)) \ge 1 - \alpha$) and reduces compliance officer triage queues by over 99.8\% with an ultra-low batch inference latency of 3.30\,ms.

Evaluated across 13 cross-domain financial benchmarks encompassing over 6.5 million entities and 45 million transactions, C-STGB establishes global Rank~\#1 performance across all 14 evaluated models. Forensic case studies demonstrate that the autonomous multi-agent swarm successfully translates geometric predictions into legally defensible FinCEN Form 111 XML filing narratives.

Future work will expand C-STGB along two high-impact axes: (i) decentralized cross-institution federated learning protocols integrated with zero-knowledge cryptographic proofs (ZK-SNARKs) for multi-jurisdictional surveillance without customer data leakage, and (ii) continual graph learning mechanisms with topological replay reservoirs to resist evolving laundering typologies in real-time.
